% Extracto editor-ready del contenido exclusivo de masas/16_hidrodinamica_y_gauge.tex.
% Las líneas 1--14 están publicadas byte-exactamente desde
% 74_rev2_realizaciones_i.tex:544--557. El propietario íntegro permanece
% preservado; este extracto conserva la continuación exclusiva.

\subsection{Comparaison des réalisations nulles électromagnétique et chromodynamique}
Le photon et les huit gluons appartiennent à des cartes nulles distinctes. Pour le photon,

\[
m_\gamma=0
\]
dans la branche électromagnétique. Pour les gluons,

\[
m_g=0
\]
dans la représentation adjointe de couleur. La nullité ne s'obtient pas en prenant la limite d'un caractère positif ; le domaine bifurque auparavant. La réalisation photonique exige Maxwell, l'invariance de jauge et deux polarisations transversales. La réalisation gluonique exige la connexion de couleur, ses huit générateurs et la restriction à l'espace physique de jauge.

La propagation bidirectionnelle peut s'écrire

\[
c_\pm=\frac{c}{1\pm\epsilon},
\qquad
\epsilon=q_-^{90}-q_+^{90},
\]
avec pour moyenne harmonique

\[
\frac{2}{1/c_++1/c_-}=c.
\]
La vitesse bidirectionnelle conserve \(c\) ; l'orientation unidirectionnelle conserve le biais. Un gap supplémentaire de l'enregistrement généalogique peut se réaliser par Proca--Stückelberg, mais ne doit pas être identifié automatiquement à \(\epsilon\) ni au dédoublement de l'action.

\Needspace{7\baselineskip}
